
\documentclass{article}
\usepackage{colortbl}
\usepackage{makecell}
\usepackage{multirow}
\usepackage{supertabular}

\begin{document}

\newcounter{utterance}

\centering \large Interaction Transcript for game `cladder', experiment `full\_v1.5\_default', episode 7433 with selfrefine{-}llama{-}3.1{-}8b{-}instruct.
\vspace{24pt}

{ \footnotesize  \setcounter{utterance}{1}
\setlength{\tabcolsep}{0pt}
\begin{supertabular}{c@{$\;$}|p{.15\linewidth}@{}p{.15\linewidth}p{.15\linewidth}p{.15\linewidth}p{.15\linewidth}p{.15\linewidth}}
    \# & \multicolumn{2}{c}{Player} && \multicolumn{2}{c}{Game Master} \\
    \hline

    \theutterance \stepcounter{utterance}  
    & & & \multicolumn{4}{p{0.6\linewidth}}{
        \cellcolor[rgb]{0.9,0.9,0.9}{
            \makecell[{{p{\linewidth}}}]{
                \texttt{\tiny{[P1$\langle$GM]}}
                \texttt{You are an expert in causal inference. The following question is not a typical commonsense query, but rather a meticulously designed question created by a professor specializing in causal inference, intended to assess the students' mastery of the course content.} \\
\texttt{Imagine a self{-}contained, hypothetical world with only the following conditions, and without any unmentioned factors or causal relationships: Vubr has a direct effect on zuph and uvzi. Zuph has a direct effect on wibl. Wibl has a direct effect on uvzi. Vubr is unobserved. The overall probability of zuph is 32\%. The probability of not zuph and uvzi is 61\%. The probability of zuph and uvzi is 29\%. Is the chance of uvzi smaller when observing zuph?} \\
\texttt{Start your answer with 'Yes' or 'No', followed by additional reasoning or evidence to support your explanation.} \\
            }
        }
    }
     \\ \\

    \theutterance \stepcounter{utterance}  
    & & \multicolumn{4}{p{0.6\linewidth}}{
        \cellcolor[rgb]{0.95,0.95,0.95}{
            \makecell[{{p{\linewidth}}}]{
                \texttt{\tiny{[P1$\leftrightarrow$P]}}
                \texttt{{[}DRAFT{]}} \\
\texttt{No.} \\
\\ 
\texttt{To answer this question, we need to apply the principles of causal inference, specifically the concept of conditional probability and the backdoor criterion. } \\
\\ 
\texttt{First, let's establish the causal relationships:} \\
\\ 
\texttt{{-} Vubr → zuph} \\
\texttt{{-} Vubr → uvzi} \\
\texttt{{-} zuph → wibl} \\
\texttt{{-} wibl → uvzi} \\
\\ 
\texttt{We are interested in the effect of observing zuph on the probability of uvzi. To do this, we need to consider the backdoor paths from zuph to uvzi.} \\
\\ 
\texttt{There are two backdoor paths:} \\
\\ 
\texttt{1. zuph → wibl → uvzi} \\
\texttt{2. zuph → uvzi (direct effect)} \\
\\ 
\texttt{However, we are given that the probability of not zuph and uvzi is 61\%, and the probability of zuph and uvzi is 29\%. This implies that there is a positive association between zuph and uvzi.} \\
\\ 
\texttt{To determine if observing zuph affects the probability of uvzi, we need to consider the effect of zuph on uvzi, controlling for the common cause Vubr. Since Vubr is unobserved, we cannot directly control for it. However, we can use the given probabilities to infer the effect of zuph on uvzi.} \\
\\ 
\texttt{Let's denote the probability of uvzi given zuph as P(uvzi|zuph) and the probability of uvzi given not zuph as P(uvzi|\textasciitilde{}zuph). We are given that P(\textasciitilde{}zuph and uvzi) = 0.61 and P(zuph and uvzi) = 0.29.} \\
\\ 
\texttt{Using the law of total probability, we can write:} \\
\\ 
\texttt{P(uvzi) = P(uvzi|\textasciitilde{}zuph)P(\textasciitilde{}zuph) + P(uvzi|zuph)P(zuph)} \\
\\ 
\texttt{Substituting the given values, we get:} \\
\\ 
\texttt{0.29 = 0.61P(uvzi|\textasciitilde{}zuph) + P(uvzi|zuph)(0.32)} \\
\\ 
\texttt{Simplifying the equation, we get:} \\
\\ 
\texttt{P(uvzi|zuph) = (0.29 {-} 0.61P(uvzi|\textasciitilde{}zuph)) / 0.32} \\
\\ 
\texttt{Since P(uvzi|\textasciitilde{}zuph) is the probability of uvzi when zuph is not observed, and we are interested in the effect of observing zuph on the probability of uvzi, we need to compare P(uvzi|zuph) to P(uvzi|\textasciitilde{}zuph).} \\
\\ 
\texttt{However, we cannot directly compare these probabilities without knowing the value of P(uvzi|\textasciitilde{}zuph). Nevertheless, we can use the given information to infer the relationship between P(uvzi|zuph) and P(uvzi|\textasciitilde{}zuph} \\
            }
        }
    }
    & \\ \\

    \theutterance \stepcounter{utterance}  
    & & \multicolumn{4}{p{0.6\linewidth}}{
        \cellcolor[rgb]{0.95,0.95,0.95}{
            \makecell[{{p{\linewidth}}}]{
                \texttt{\tiny{[P1$\leftrightarrow$P]}}
                \texttt{{[}SELF{-}FEEDBACK{]}} \\
\texttt{FORMAT: ok} \\
\texttt{RULES: ok} \\
\texttt{STRATEGY: problem} \\
\texttt{REASON: The answer does not directly address the question of whether the chance of uvzi is smaller when observing zuph, but instead provides a detailed analysis of the causal relationships and conditional probabilities, which is not a clear yes or no answer.} \\
            }
        }
    }
    & \\ \\

    \theutterance \stepcounter{utterance}  
    & \multicolumn{4}{p{0.6\linewidth}}{
        \cellcolor[rgb]{0.561,0.737,0.561}{
            \makecell[{{p{\linewidth}}}]{
                \texttt{\tiny{[P1$\rangle$GM]}}
                \texttt{No.} \\
\\ 
\texttt{The probability of uvzi given zuph, P(uvzi|zuph), is not directly calculable from the given information. However, we can infer that observing zuph does not necessarily make the chance of uvzi smaller. The positive association between zuph and uvzi, as indicated by the given probabilities, suggests that observing zuph may actually increase the chance of uvzi.} \\
            }
        }
    }
    & & \\ \\

    \theutterance \stepcounter{utterance}  
    & & & \multicolumn{2}{p{0.3\linewidth}}{
        \cellcolor[rgb]{0.9,0.9,0.9}{
            \makecell[{{p{\linewidth}}}]{
                \texttt{\tiny{[GM$|$GM]}}
                \texttt{No} \\
            }
        }
    }
    & & \\ \\

    \theutterance \stepcounter{utterance}  
    & & & \multicolumn{2}{p{0.3\linewidth}}{
        \cellcolor[rgb]{0.9,0.9,0.9}{
            \makecell[{{p{\linewidth}}}]{
                \texttt{\tiny{[GM$|$GM]}}
                \texttt{no} \\
            }
        }
    }
    & & \\ \\

    \theutterance \stepcounter{utterance}  
    & & & \multicolumn{2}{p{0.3\linewidth}}{
        \cellcolor[rgb]{0.9,0.9,0.9}{
            \makecell[{{p{\linewidth}}}]{
                \texttt{\tiny{[GM$|$GM]}}
                \texttt{game\_result = WIN} \\
            }
        }
    }
    & & \\ \\

\end{supertabular}
}

\end{document}
